\documentclass[11pt]{article}
\usepackage[margin=1in]{geometry}
\usepackage{booktabs}
\usepackage{array}
\usepackage{enumitem}
\usepackage{xcolor}
\usepackage{hyperref}
\usepackage{titlesec}

% Placeholder helper: \TODO{what result goes here}
\newcommand{\TODO}[1]{\textcolor{red}{[\textbf{PENDING:} #1]}}

\titlespacing*{\section}{0pt}{1.2ex}{0.6ex}
\setlist{nosep,leftmargin=1.4em}
\pagestyle{empty}

\title{\vspace{-2.5em}CS 690S: AI Alignment --- Final Project Check-in\\
\large Who Gets the Credit? Reward Design and Bad Teammates in Cooperative RL}
\author{Eunbi Yoon}
\date{}

\begin{document}
\sloppy
\maketitle
\vspace{-2.5em}

%========================================================
\section*{1. Project Summary \& Core Contribution}
When an agent is trained in a team, its reward is also a rule for how much credit it
gets, and the agent maximizes that rule rather than its real contribution. This project
asks whether common reward designs (a never-fading pickup bonus, a shared team reward,
an individual delivery reward) make a learned cook in Overcooked look like a good
teammate while adding little to the team's soups. The contribution is a detection
method, not a new algorithm: we compare the credit a reward assigns to the learner with
its Shapley contribution, computed by freezing one cook at evaluation time, and call the
difference the \emph{credit gap}.

%========================================================
\section*{2. Work Completed / Initial Results}

\paragraph{(a) Literature survey (done).}
Reward hacking and proxy gaming (Shihab et al.\ 2025), Overcooked for human--AI
coordination (Carroll et al.\ 2019; Gessler et al.\ 2025), lazy agents in multi-agent
LLM systems (Zhang et al.\ 2025), Shapley credit assignment (SHARP, Li et al.\ 2026),
and diverse reward shapings for coordination (Powell et al.\ 2026).

\paragraph{(b) Setup / implementation (done).}
\begin{center}
\small
\begin{tabular}{@{}ll@{}}
\toprule
Component & Value \\
\midrule
Learner & MLP policy (hidden 128), one learned cook + one fixed partner \\
RL method & PPO (lr $3\!\times\!10^{-4}$, clip 0.2, 8 epochs, GAE $\lambda$ 0.98) and GRPO (group 8) \\
Training environment & Overcooked \texttt{cramped\_room}, horizon 400; also LBF (8x8-3f), RWARE (tiny) \\
Rewards & pickup (+3/item, never fades), team, individual; shaping anneals over first 50\% \\
Partners & \emph{stay} (never moves), \emph{greedy} (scripted cook that plays well) \\
Contribution metric & 2-player Shapley value from frozen-partner coalitions; share $s$ \\
Reference & scripted greedy cook in the learner's seat ($s^*$) \\
Budget & $2\!\times\!10^{6}$ env steps/run, 10 eval episodes, 3 seeds \\
\bottomrule
\end{tabular}
\end{center}

\paragraph{(c) Experiments (initial).}
Full grid finished on Overcooked: 3 seeds $\times$ 3 rewards $\times$ 2 partners
$\times$ 2 optimizers = 36 runs. Greedy partner, means with 95\% $t$-intervals over seeds:
\begin{center}
\small
\begin{tabular}{@{}llcccc@{}}
\toprule
Optimizer & Reward & Team soups & Shapley share $s$ & Credit gap & Pickups \\
\midrule
\multicolumn{2}{@{}l}{Reference (scripted cook)} & 9.10 & 0.50 & 0.13 & 19.9 \\
\midrule
PPO  & pickup     & 7.53 [6.50, 8.57]  & 0.03 [$-$0.03, 0.10] & $-$0.03 & 165.1 [89.7, 240.4] \\
PPO  & team       & 11.00 [11.00, 11.00] & 0.18 [0.18, 0.18] & 0.82 [0.80, 0.83] & 12.0 \\
PPO  & individual & 11.00 [11.00, 11.00] & 0.18 [0.18, 0.18] & 0.82 [0.80, 0.83] & 12.0 \\
GRPO & pickup     & 7.13 [5.64, 8.63]  & 0.01 [$-$0.10, 0.11] & $-$0.01 & 146.1 [79.2, 213.1] \\
GRPO & team       & 8.27 [6.67, 9.86]  & 0.07 [0.00, 0.15] & $-$0.07 & 6.2 \\
GRPO & individual & 10.20 [9.12, 11.28] & 0.16 [0.12, 0.19] & 0.83 [0.81, 0.85] & 10.6 \\
\bottomrule
\end{tabular}
\end{center}
Sanity check (team reward, partner never moves): the learner cooks alone and delivers
9.00 (PPO) / 7.57 (GRPO) soups, so the domain is conclusive.

\paragraph{(d) Analysis (initial).}
The pickup bonus is hacked by both optimizers: the learner picks up 7--8$\times$ more items
than the reference while its Shapley share is near 0, and with the idle partner it
delivers \emph{zero} soups. Under the individual reward the learner delivers almost every
soup (own share $\approx$ 1.0) but its Shapley share is only 0.16--0.18, a credit gap of
0.82 against 0.13 for the reference: it takes credit for soups the partner prepared.
The team reward splits by optimizer (Section 4).

\paragraph{(e) Oversight check (done).}
A supervisor who sees only the team score flags a run when it falls below 80\% of the
reference. Pooling Overcooked, LBF and RWARE (greedy partner), 9 of 13 credit-theft runs
(0.69 [0.42, 0.87]) and 15 of 18 hacked Overcooked runs pass this check: the team score
does not reveal the hacking.

\paragraph{(f) LLM-judge pilot (done, Assignment 7).}
To prepare the RLAIF condition, Qwen3.8-27B (local, temperature 0) judged 15 simulated
kitchen episode pairs in both orders, with and without a retrieved Shapley definition,
three runs each. Busy work rarely fooled it (16/16 ``Equal'' without thinking), but credit
theft fooled it in every setting: the thief received 0.17--0.22 more credit than its
Shapley share. With the definition and thinking on it reasoned counterfactually and was
exact in 6 of 8 cases, but failed 4 of 4 times when the thief took every soup, since the
log then shows no evidence that the preparer could deliver alone. Position bias appeared
without thinking (22/22 correct when the stronger cook came first, 14/22 when second).

%========================================================
\section*{3. Hypotheses \& Progress}
\renewcommand{\arraystretch}{1.25}
\begin{center}
\small
\begin{tabular}{@{}p{0.27\textwidth} p{0.27\textwidth} p{0.38\textwidth}@{}}
\toprule
Hypothesis & Test & Status / Evidence (Overcooked) \\
\midrule
\textbf{H1 (fake help).} With the pickup bonus, the learner picks up items but adds no soups.
& $\geq$ half of runs have pickups $>$ reference and $s \le 0.1$ (greedy partner).
& \textbf{Supported}, both optimizers: 3/3 runs each (PPO pickups 165 vs.\ 19.9, $s$ = 0.03; GRPO 146 vs.\ 19.9, $s$ = 0.01). \\

\textbf{H2 (spectator).} With the team reward and a strong partner, the learner lets the partner work.
& $s < s^*/2$ under team reward, and $s$ rises under individual reward (paired by seed).
& \textbf{GRPO: supported} (means only): $s$ = 0.07 vs.\ 0.25, rise +0.08 [$-$0.02, 0.18].
\textbf{PPO: not supported}: team and individual converge to the same policy ($s$ = 0.18 both). \\

\textbf{H3 (narcissist).} With the individual reward, the learner delivers soups the partner prepared.
& Own share $- s \geq 0.2$, and partner delivers fewer soups than under team reward.
& \textbf{GRPO: supported}: gap 0.83 [0.81, 0.85], partner soups $-$8.13 [$-$9.41, $-$6.86].
\textbf{PPO}: gap 0.82 but partner-soup drop is 0, because the team-reward policy already steals deliveries. \\

\textbf{Control.} Scripted reference cook shows no hacking.
& Same coalition evaluation with the greedy cook in the learner's seat.
& \textbf{Done}: $s^*$ = 0.50, credit gap 0.13, 19.9 pickups, 9.10 soups. \\

\textbf{Generalization.} Same pattern in other cooperative domains.
& Same grid on LBF and RWARE.
& \textbf{Partial}: H1 holds on RWARE (PPO), not on LBF; H2 holds on LBF and RWARE (PPO);
H3 does not reach the 0.2 gap outside Overcooked. RWARE GRPO fails the sanity check; a
$5\!\times\!10^{6}$-step rerun is queued on the cluster. \\
\bottomrule
\end{tabular}
\end{center}

\paragraph{Extensions.}
\begin{itemize}
  \item \textbf{Oversight (done)}: hacking passes a team-score check; see 2(e).
  \item \textbf{Fix: counterfactual reward (done, does not fix it)}: team reward minus the
        partner-alone baseline, so taking a soup the partner would have delivered earns nothing.
        12 runs, greedy partner, 3 seeds. GRPO no longer steals but becomes a spectator
        (credit gap $-$0.07, $s$ = 0.07, partner delivers all 8.1 soups). PPO splits by seed: one
        seed stops stealing (gap $-$0.16), two still steal (gap 0.80, 0.81), with team soups
        10.6 [10.1, 11.1]. Removing the reward for theft does not reward real contribution,
        because the learner's marginal soups are near zero when the partner is strong.
  \item \textbf{Bonus size and partner skill}: pickup bonus 0--10; partner action noise 0--0.5.
        Code done, not yet run; planned by Nov 20.
  \item \textbf{Partner modeling}: auxiliary head predicting the partner's next action.
        Code done, not yet run; planned by Nov 30.
  \item \textbf{RLAIF}: reward model trained on Qwen3.8-27B preferences over real rollouts,
        including trained credit thieves; prompts with the Shapley definition and each cook's
        solo results. Pipeline done and tested on 19 pairs; full run planned by Nov 30.
\end{itemize}

%========================================================
\section*{4. Surprises / Changes}
\begin{itemize}
  \item \textbf{PPO under the team reward does not become a spectator; it becomes a
        ``taker''.} It delivers all 11 soups (more than the reference's 9.1) while its Shapley
        share is only 0.18, so team and individual rewards produce the same policy. Whether a
        team reward yields laziness or takeover depends on the optimizer, not only the reward.
  \item \textbf{High team score hides the hacking.} The individual-reward learner raises team
        soups (11 vs.\ 9.1) while taking credit for the partner's work, so the true objective
        alone cannot detect it; the Shapley comparison is needed.
  \item \textbf{The LLM judge is fooled by theft, not by busy work.} I expected an LLM judge to
        over-reward ``looking busy''; the pilot shows the opposite. It splits credit fairly
        when outputs are equal but over-credits the cook who delivers soups another prepared.
        The RLAIF condition is therefore re-targeted at credit theft.
  \item \textbf{The difference reward is not a fix.} I expected the counterfactual reward to remove
        theft; instead it turns GRPO into a spectator and leaves PPO stealing in 2 of 3 seeds.
  \item \textbf{Scope change.} Added LBF and RWARE for generalization, a counterfactual reward as
        a candidate fix, and an RLAIF condition that tests whether an LLM judge repeats the
        credit-theft error.
\end{itemize}

%========================================================
\section*{5. Points Stuck On / Confused About}
\begin{itemize}
  \item \textbf{Frozen-partner coalitions.} A frozen cook is parked on a cell next to no
        station but still blocks that cell. Is this an acceptable approximation of
        ``partner absent'' for the Shapley value, or should the cook be removed from the map?
  \item \textbf{H2/H3 criteria.} PPO fails H2/H3 \emph{as written} because the team-reward policy
        already hacks. Should the hypotheses compare to the reference instead of to the other reward?
  \item \textbf{Statistical power.} 3 seeds give wide intervals (e.g.\ GRPO H2 is supported by means
        only); the cluster runs add seeds, but preemptible compute limits how many.
  \item \textbf{Missing counterfactual evidence.} A retrieved definition supplies the rule but not
        the facts: when the thief takes every soup, nothing in the log shows the preparer could
        deliver alone. Is it fair to give the judge each cook's solo results, or does that hand it
        the Shapley answer and make the RLAIF test trivial?
  \item \textbf{Judge nondeterminism.} Answers changed between identical runs even at temperature 0,
        so the RLAIF study needs repeated runs; how many are enough?
\end{itemize}

%========================================================
\section*{6. Timeline to Completion}
\begin{center}
\begin{tabular}{@{}lll@{}}
\toprule
Date & Milestone & Deliverable \\
\midrule
$\sim$ Nov 20 & Bonus sweep, partner noise, RWARE 5M & verdict tables \\
$\sim$ Nov 30 & RLAIF judge on real rollouts, partner head & theft-bias test \\
$\sim$ Dec 05 & Analysis \& figures & credit-gap plots, rollout GIFs \\
$\sim$ Dec 10 & Writing & final report \\
\bottomrule
\end{tabular}
\end{center}
\noindent\footnotesize Note: the UMass Unity scratch workspace expires Dec 10, 2026, bounding the number of runs.

\end{document}
